\section{Carrying The Residue}

A repeated trial is not a row of fresh beginnings. If every pass begins from
zero, the passes may resemble one another, but they do not refine one another.
They are separate questions. The grammar of alpha cannot use that kind of
repetition, because \(\alpha\) is not read from a single isolated sample. It
is read from a bounded process in which each sample inherits what the previous
sample did not settle.

The inherited state has two essential pieces: seed and residue. The seed is
the carried selector by which the next sample is placed. It tells the grammar
which admissible face of the aperture is current. It is not an unaccountable
whim and not an appeal to randomness as explanation. It is the part of the
trial state that makes the next selection repeatable as a successor of the
old one. Without a carried seed, the next sample may be possible, but it is
not yet the next sample of this trial.

The residue is the unfinished remainder. It records the fraction of a turn,
the leftover strain, or the unabsorbed reading that has not yet crossed the
trial's counting boundary. A residue that disappears between trials is not a
residue; it is a discarded error. The grammar forbids that discard. If a
finite pass leaves an unread remainder, the next pass must begin with that
remainder in the ledger.

This is why the trial advances by carrying rather than restarting. A pass
selects an aperture face, samples the finite neighborhood allowed by that
face, projects the mean remainder onto the current lattice of the trial, and
adds it to the residue already present. If the sum crosses a full turn, the
grammar counts the completed turn. If it does not, the grammar carries the
remainder. In either case, the new state is not blank. It is the old state
extended.

The wrap is not a loss of information. When a residue crosses one full turn,
the completed turn is counted and the remainder is restored as the starting
residue for the next pass. The count and the residue therefore separate two
faces of the same event. The count says what has closed. The residue says
what remains open. A trial that keeps only the count loses the local memory
needed for the next aperture. A trial that keeps only the residue loses the
closed loop that makes charge billable.

The seed is carried through the same extension. Once a face has been sampled,
the selector advances. The next face is not chosen from outside the record.
It is obtained by the successor state of the trial. The grammar may therefore
read the next sample as belonging to the same repeated process. A stable
procedure does not mean that the same face must be sampled every time. It
means that the rule by which faces are sampled remains part of the record.

This distinction protects the trial from a common statistical shortcut.
Repeated measurements can be averaged only when the alphabet of measurement
and the update rule have remained stable. If the alphabet changes between
passes, the mean is a mixture of incompatible ledgers. If the update rule
changes, the spread no longer reports the residue of one procedure. The
grammar of alpha requires a stronger condition: not merely many samples, but
many samples whose seed and residue form one bounded history.

The finite bound is equally important. A residue carried without a horizon
would become an excuse to defer every debt. The grammar must know how much
of the trial has been spent, how many passes have been permitted, and what
remainder is present when the bound is reached. A bounded trial may stop with
residue still visible. That is not failure. It is the honest statement of
what the trial has read and what it has not read.

The laser-tracking example gives a clean image. A reader who emits a pulse
and records its return has not erased the previous pulse. Each emission and
return pair constrains the path already being tracked. A long record is not
many beginnings; it is one growing chain of commitments. If the tracked body
accelerates, later returns force new constraints into the same ledger. The
chain grows because the earlier constraints remain in force.

The same form governs invisible motion between anchors. If two readers refine
an interval and no new event appears, the grammar does not say that the
interval was empty. It says that the carried residue has been driven below
the current distinguishability threshold for that trial. If a later shared
anchor exposes a disagreement in value, slope, or bending moment, the
disagreement becomes residue that must be carried forward. Repeatability is
the rule that makes such residue communicable.

For alpha, the carried residue has a special consequence. A single aperture
sample cannot decide the fine-structure reading, because the reading belongs
to the relation between loop count, strain, and value. Each pass samples one
face under one finite selection. The residue of that pass supplies the next
pass with its starting debt. The seed supplies the next pass with its
admissible selector. The trial becomes an orbit of bounded obligations.

The grammar therefore reads finite repetition in three layers. It permits a
next pass only as an extension of the prior state. It carries seed and residue
as the state that makes the next pass continuous with the last. It bounds the
sequence so the carried residue remains a readable debt rather than a promise
of future absolution. Only under those rules can repeated trials contribute
to a finite extrapolation of \(\alpha\).

The next question is what the aperture itself can read. The carried trial
does not sample a formless opening. It samples a bracket with three faces.
Those faces are the rungs on which the alpha reading will stand: charge,
mass, and value.
